\chapter{机器学习理论基础}
\label{chap:statistical-learning-framework}

我们用一批已标注的邮件训练垃圾邮件分类器，希望它也能判断新收到的邮件。训练邮件上的准确率只能检验分类器对这批输入的回答：能够利用共同规律的分类器与逐条记住训练数据的分类器，都可能在这里得到高分。要评价学习的结果，还须考察它们面对未见邮件时的错误。

把另一批邮件留作测试，可以估计当前分类器的错误率。学习算法的保证却还涉及另一层变化：重新抽取训练邮件后，算法可能选出不同的分类器。有限观测何时足以支持对未来错误率的判断，因而取决于数据怎样产生、算法怎样选择结果，以及我们允许多大的预测误差和训练失败概率。机器学习理论把这些关系写成明确条件，使一次测试的结果与重复训练的保证能够分别分析。

为分析这些条件，本章建立第~\ref{chap:statistical-learning-framework}--\ref{chap:transfer-learning-theory}章共用的学习问题框架。邮件二分类给出观测、假设和风险的具体含义，PAC标准进一步约束算法在随机训练中的成功概率。记忆型反例与没有免费的午餐定理说明经验拟合的局限，风险分解和模型选择则解释限制假设空间的收益与代价。研究历史及当前前沿把这一框架放到更广的理论问题中，说明标签、损失、计算和环境变化为何需要分别讨论。

\section{学习问题}
\label{sec:learning-problem}
\label{sec:statistical-learning-problem}

在邮件分类中，一条训练记录包含邮件与标签，算法交付一个分类器，未来错误率用来评价它。在回归或聚类中，这些对象会改变；形式相同的风险界也就可能约束不同的结果。我们先把观测、候选决策、算法输出、计分方式与随机来源分开规定，再用二分类逐一具体化。这些职责由第~\ref{chap:statistical-learning-framework}--\ref{chap:transfer-learning-theory}章共享，对象的具体形式则随任务确定。

\subsection{基本设定}

先看算法收到的一条记录。邮件分类记录邮件及其类别，回归记录输入与数值响应，聚类记录待分组的对象。这些允许记录组成\emph{观测空间}（observation space），记为$Z$；损失函数以其中的一条观测评价学习结果。指定$Z$时也须确定记录的单位：单条记录、相互依赖的序列和整块环境数据可能包含不同信息，不能直接按同一种独立样本计数。

观测确定以后，还要给学习结果选定比较对象。全部阈值分类器、受限参数的回归函数或固定中心数的聚类方案，都可以组成一组候选决策，称为\emph{比较决策空间}（comparison decision space），记为$\mathcal W$。最低风险及超额风险相对于这组候选计算。另以\emph{算法输出空间}（algorithm output space）$\mathcal W_{\mathrm{out}}$记录算法能够交付的结果：直接选择候选时两者相同，组合候选或交付其他预测器时，输出范围可以更大。

评价决策还需要一条计分规则。\term{损失函数}（loss function）把决策与观测映射为代价，较小的损失表示较好的评价结果。至于算法会看到哪些观测、观测之间是否独立、未来损失对哪个总体取平均，则由数据来源及抽样机制规定。样本之外的标签、环境信息和辅助变量是否可见，也属于算法输入的一部分。下面将这些对象合在同一个形式设定中，分类与独立同分布只是它的具体取法。

\begin{definition}[学习问题的形式设定]
\label{def:learning-problem-setting}
一个学习问题至少需要规定观测空间$Z$、比较决策空间$\mathcal W$、算法输出空间
$\mathcal W_{\mathrm{out}}$、损失函数、允许的数据来源、抽样机制和算法可见的信息。

在$Z$与$\mathcal W^+=\mathcal W\cup\mathcal W_{\mathrm{out}}$上约定相应的可测结构，
并要求损失
\[
 \ell:\mathcal W^+\times Z\to\mathbb R
\]
联合可测。

抽样机制规定算法所见观测$S$的联合分布。内部随机数$U$定义在另一个概率空间并独立于
$S$；学习算法$A$是从允许输入到$\mathcal W_{\mathrm{out}}$的可测映射，输出
$A(S,U)$。确定性算法省略$U$。
\end{definition}

在这个设定里，$Z$与$\ell$共同说明一条记录怎样为决策计分。$\mathcal W$提供比较基准，$\mathcal W_{\mathrm{out}}$约束交付结果；因此，评价时比较哪些决策与训练时可以返回哪些决策，是两项分别指定的条件。

允许的数据来源限定保证覆盖的生成机制，通常不会作为已知信息交给算法。抽样机制进一步规定训练观测彼此之间、训练观测与评价观测之间的关系。独立随机数$U$记录算法自身的随机选择，$A$则根据可见输入与$U$产生决策。如此才能区分数据中的不确定性与算法引入的不确定性。

先取最常用的同分布基线。以$\mathfrak P$收集允许来源，给定$\mathcal P\in\mathfrak P$后，从中独立抽取$n$条观测：
\[
 S=(z_1,\ldots,z_n)\sim\mathcal P^n.
\]
二分类、多分类与一般损失章节以此为起点。迁移学习另行规定源域、目标域、各样本量或环境层抽样，相应风险与概率也随之改变。各章沿用定义~\ref{def:learning-problem-setting}的对象职责，抽样关系则须在当地重新指定。

应用这一设定时，样本量必须随观测单位解释。独立记录的条数、序列的长度与环境块的个数，不能未经分析就相互替换。结论还须注明允许哪些分布：覆盖全部相关分布、只覆盖受限分布族或覆盖特定跨域关系，对学习算法提出的要求不同。

同样，风险与容量必须随比较空间解释。“误差小”可以指总风险小，也可以指相对于$\mathcal W$中最低风险的差距小。若算法始终返回$\mathcal W$中的决策，称为\term{恰当学习}（proper learning）；若允许返回$\mathcal W$之外而在$\mathcal W_{\mathrm{out}}$中的决策，称为\term{非恰当学习}（improper learning）。两种学习仍以$\mathcal W$的最低风险比较，改变的是可交付结果及相应的算法接口。

保证中的概率也要明确对谁取平均。新观测决定固定决策的期望损失，训练样本与算法随机数决定训练成功的概率，环境抽样还可能增加一层随机性；相应的期望、高概率和逐环境保证约束不同对象。可测性与可分性使这些概率及经验过程上确界有定义。统计保证控制样本量，可计算性要求程序停机，高效性还限制训练与预测时间。分类、实值损失和跨域学习可以复用这些职责，但改变任何一项都须重新核对保证，并在具体定理中补齐有界性、凸性或覆盖条件。

\paragraph{可测性与可分性约定}

算法看过样本再选决策，风险分析因此常须同时比较一组损失。对需要取上确界的$\mathcal W'\subseteq\mathcal W^+$，先将其损失写成函数类：
\[
 \mathcal F_\ell(\mathcal W')
 =\{z\mapsto\ell(w,z):w\in\mathcal W'\}.
\]
直接给定函数类时仍记为$\mathcal F$。本部分要求一致有界类中的函数均可测，且有可数子类$\mathcal F_0$，使任意$f\in\mathcal F$都是其中一个序列的逐点极限。若$\mathcal F=\mathcal F_\ell(\mathcal W')$，为$\mathcal F_0$的每个函数固定一个代表决策，得到可数$\mathcal W_0\subseteq\mathcal W'$及$\mathcal F_0=\mathcal F_\ell(\mathcal W_0)$。学习、压缩、重构及其他数据依赖输出映射还须可测。

这些条件使不可数比较可以回到可数对象上。逐点逼近先保证固定有限样本的经验平均收敛，一致有界再通过有界收敛保证期望平均收敛，相关上确界因此可在可数子类上计算。给定$\rho>0$，按$\mathcal W_0$的固定枚举选第一个满足$\hat R_S(w)\leq\inf_{v\in\mathcal W_0}\hat R_S(v)+\rho$的决策，就得到可测选择：下确界和比较事件都可测，而逐点逼近给出$\inf_{v\in\mathcal W_0}\hat R_S(v)=\inf_{v\in\mathcal W'}\hat R_S(v)$。这也是原集合上的近似ERM。无法建立可数逼近或可测选择时，须改用外概率或另加经验过程条件；无界损失则另行检查可积性与矩条件。

\subsection{二分类例子}

现在取邮件作为观测，把统一设定落实到二分类。令$Z=X\times Y$、$Y=\{0,1\}$，以假设空间$\mathcal H$同时作为比较空间与输出空间，并用$\ell(h,(x,y))=\ell_{01}(h(x),y)$记录预测是否出错。接下来的阈值例子将说明有限标签怎样限制候选假设，也为可实现与不可知两种比较基准提供共同对象。

算法用已知类别的邮件选择预测规则。标签即使完全正确，也只约束已经观察到的邮件，可能留下几条对训练数据回答相同、对新邮件回答不同的规则。

设每封邮件都已有一个$[0,1]$中的可疑程度分数。较高分数表示评分程序将这封邮件判断为更可疑，并不必然表示相应的垃圾邮件概率。我们固定评分程序，只学习分类的分界：分数达到分界时判为垃圾邮件，低于分界时判为正常邮件。这个分界值称为\term{阈值}（threshold），学习任务就是根据标签选择它。

用四封人为构造的邮件检验这一选择。分数0.1与0.3对应正常邮件，0.7与0.9对应垃圾邮件。阈值0.4和0.6都把前两封判为正常、后两封判为垃圾，所以训练错误都为零。两个阈值之间没有训练邮件，已见标签便无法进一步判断阈值应落在哪里。

分数为0.5的新邮件正好落在两条规则的分歧处。由于$0.5\geq0.4$，第一条规则预测垃圾邮件；由于$0.5<0.6$，第二条规则预测正常邮件。表~\ref{tab:threshold-email-decisions}把已见与未见输入放在一起比较。新邮件尚无真实标签，我们不能判定哪条规则正确，但能据此确认：四封训练邮件未能排除这项分歧。

\begin{table}[htbp]
 \centering
 \begin{tabular}{cccc}
  \toprule
  邮件分数 & 已知类别 & 阈值0.4的预测 & 阈值0.6的预测\\
  \midrule
  0.1 & 正常 & 正常 & 正常\\
  0.3 & 正常 & 正常 & 正常\\
  0.7 & 垃圾 & 垃圾 & 垃圾\\
  0.9 & 垃圾 & 垃圾 & 垃圾\\
  \midrule
  0.5 & 尚未标注 & 垃圾 & 正常\\
  \bottomrule
 \end{tabular}
 \caption{四封构造的训练邮件与一封未标注邮件上的阈值预测。两个阈值分对全部训练记录，却在分数0.5处产生分歧；末行没有真实标签，故不能据它判断谁正确。}
 \label{tab:threshold-email-decisions}
\end{table}

在这个例子中，邮件分数是分类器可读取的\term{输入}，记为$x$；人工给出的真实类别是\term{标签}，记为$y$。约定垃圾邮件为1、正常邮件为0，所有允许输入组成$X=[0,1]$，所有标签组成$Y=\{0,1\}$。实际输入也可以是多个特征组成的向量；输入与标签的职责不随这种表示变化。

阈值一旦选定，就得到一条对每个允许输入给出类别的规则，称为\term{假设}（hypothesis），写作函数$h:X\to Y$。例如$h(0.5)=1$只陈述这条规则对分数0.5的预测；真实标签仍须另行观察，不能由预测值替代。

算法允许选择的全部假设组成非空的\term{假设空间}（hypothesis space），记为$\mathcal H$。当前只调整阈值，故以全部允许$t$为索引写出候选：
\begin{equation}
 \mathcal H_{\mathrm{thr}}=\{h_t:t\in[0,1]\},\qquad
 h_t(x)=\begin{cases}
  1,&x\geq t,\\
  0,&x<t.
 \end{cases}
 \label{eq:threshold-hypothesis-class}
\end{equation}
下标$t$标明采用哪条规则。$h_{0.4}$与$h_{0.6}$是其中两个成员，训练选择$t$之后，预测再使用对应的$h_t$。

阈值空间把不同输入的预测联系起来：越过阈值后，分数继续上升不会使预测变回正常邮件。训练数据并没有迫使我们采用这条限制，逐条记住四个输入并任意回答其他输入，也可以分对训练邮件。预先限定候选假设或规定选择偏好，称为\term{归纳偏置}（inductive bias）。它补充了样本之外的规则，使算法能在未见输入上延伸已有判断。

奥卡姆剃刀表达优先采用简单解释的想法。在阈值例子里，描述一个分界值即可指定假设，记忆型规则则须保存每个训练输入。不过，描述长度取决于编码方式，不能仅凭“简单”二字推出泛化。还需要证明这种表示限制怎样控制统计误差。

\paragraph{损失、期望风险与经验风险}

比较候选时，先把预测错了多少转成代价。\term{损失}（loss）为一次预测计分；当前对两种误判同等计分，错分记1，正确记0。这称为\term{零一损失}（zero-one loss）。以$h(x)$为预测、$y$为真实标签，有
\begin{equation}
 \ell(h,(x,y))=\ell_{01}(h(x),y)=\begin{cases}
  1,&h(x)\neq y,\\
  0,&h(x)=y.
 \end{cases}
 \label{eq:zero-one-loss}
\end{equation}
若误删正常邮件的代价更大，这一对称计分就须更换。下文先固定零一损失，以免比较时改变评价口径。

一次预测的损失只评价一封邮件。若分类器分错100封已标注邮件中的3封，样本平均损失是$3/100=0.03$；它直接可算，却只描述这100封邮件。我们要评价未来表现，还须计算随机到来的一封邮件被分错的概率，这需要知道不同邮件及其标签出现的频率。

以联合分布$\mathcal P$表示这种来源。它同时记录输入分数出现的频率与相应标签的概率，作用于$z=(x,y)\in Z=X\times Y$。例如，两条阈值规则只在很窄的区间上不同，若邮件恰好大量集中在那里，二者的未来错误率仍可相差很多。算法通常只看到来自$\mathcal P$的样本，并不知道完整分布。

按$\mathcal P$平均未来损失，就得到$h$的期望风险。零一损失只取0或1，其平均也就是分类错误的概率：
\begin{equation}
 R_{\mathcal P}(h)
 =\mathbb E_{(x,y)\sim\mathcal P}[\ell_{01}(h(x),y)]
 =\mathbb P_{(x,y)\sim\mathcal P}(h(x)\neq y).
 \label{eq:population-classification-risk}
\end{equation}

算法无法直接计算这个未知分布平均，但可计算样本平均。对$S=(z_1,\ldots,z_n)$、$z_i=(x_i,y_i)$，定义经验风险为
\begin{equation}
 \hat R_S(h)=\frac1n\sum_{i=1}^n\ell_{01}(h(x_i),y_i).
 \label{eq:empirical-classification-risk}
\end{equation}
分子累计样本错误数，除以$n$后给出这份样本中的错误比例。

下标区分了两种平均：$R_{\mathcal P}$依赖总体分布，$\hat R_S$依赖已经观察到的样本。所谓“经验”就是依据样本计算。$S$取训练集时，它给出训练平均损失；换成独立测试集时，它给出另一次对未来损失的样本估计。两者都不会因为能够直接计算而等于未知的总体风险。

我们还可以在同一个假设上比较总体平均与样本平均。两者的差$R_{\mathcal P}(h)-\hat R_S(h)$称为\term{泛化差}（generalization gap），可能为正，也可能为负。它考察评价样本造成的偏离。后面的超额期望风险则比较学习结果与基准的两个总体风险，考察学习结果距基准有多远。

四封训练邮件只让我们知道两条阈值的经验风险都为零，尚不足以计算总体风险。为做一次可复算的教学比较，另设分数在$[0,1]$上均匀分布，且真实垃圾邮件恰为分数不低于0.6的邮件。于是$h_{0.6}$处处正确，$h_{0.4}$只在$[0.4,0.6)$误把正常邮件判为垃圾。该区间概率为0.2，两条规则的期望风险便分别为0和0.2。结果来自补充的分布与标签规律，不能归为四封训练邮件本身提供的证据。

\paragraph{样本与学习算法}

用样本比例推断未来风险，还须连接两批数据的来源。\term{独立同分布}（independent and identically distributed, i.i.d.）要求训练观测彼此独立，且各自服从评价所用的同一个$\mathcal P$。$n$个观测的联合分布因此是
\begin{equation}
 S=\bigl((x_1,y_1),\ldots,(x_n,y_n)\bigr)\sim\mathcal P^n.
 \label{eq:training-sample-distribution}
\end{equation}
这个理想化条件需要针对数据检查。重复转发可能造成依赖，部署后的邮件类型变化可能使评价不同于训练。后一情形须加入第~\ref{chap:transfer-learning-theory}章分析的跨域条件。

从样本到假设的选择过程称为\term{学习算法}（learning algorithm），记为$A$；一次训练输出$\hat h=A(S)$。例如，算法可以比较各阈值的训练错误数并选取最小者。若多个阈值并列，还须有确定的选解约定。$A$规定怎样选择，$\hat h$是这份样本下选出的结果。

本章先令$\mathcal W_{\mathrm{out}}=\mathcal H$，采用第~\ref{sec:learning-problem}节定义的恰当学习。若候选只有单阈值规则，算法却组合规则得到不再属于该空间的分类器，就须将它纳入更大的输出空间；后文使用非恰当输出时会明示这种变化。

样本变化会使同一算法选出不同假设。一批邮件可能包含许多接近真实分界的输入，另一批恰好缺少它们，训练阈值便可能不同。因此，固定一个训练结果估计其风险，与重复抽样考察算法有多大机会得到低风险结果，是两层评价。

图~\ref{fig:learning-two-randomness}将两层随机性分开画出。固定$\hat h$后，$R_{\mathcal P}(\hat h)$只对新邮件平均；重新抽取$S$时，$A(S)$改变，它的总体风险也随之改变。学习保证控制的是后一次变化产生高风险结果的概率。算法若使用内部随机数，也须把这些随机选择纳入成功概率。

\begin{figure*}[htbp]
  \centering
  \includegraphics[width=164.600mm]{figures/03-foundations/02-learning-theory/generated/learning-theory-two-randomness/figure.png}
  \caption{两层随机性的评价对象。固定训练输出后，右侧新观测的平均损失给出期望风险；左侧重新抽样则改变输出，学习保证控制这种变化造成高风险结果的概率。}
 \label{fig:learning-two-randomness}
\end{figure*}

\subsection{一般学习的风险与可学习性}

式~\eqref{eq:population-classification-risk}与\eqref{eq:empirical-classification-risk}已经给出了分类中的两种平均。对回归函数、聚类中心等决策，可沿用相同计算关系：以一般观测$z$替代输入标签对，以决策$w$替代分类器$h$，损失仍由任务指定。这样推广的是风险的定义形式，零一损失下“错误概率”的解释则不随之推广。

\begin{definition}[一般学习问题的期望风险]
\label{def:general-population-risk}
设$\mathcal P\in\mathfrak P$。对任意$w\in\mathcal W^+$，若$\ell(w,\cdot)$在$\mathcal P$下可积，则$w$的期望风险为
\[
 R_{\mathcal P}(w)=\mathbb E_{z\sim\mathcal P}\ell(w,z).
\]
\end{definition}

\begin{definition}[一般学习问题的经验风险]
\label{def:general-empirical-risk}
给定样本$S=(z_1,\ldots,z_n)\in Z^n$，任意$w\in\mathcal W\cup\mathcal W_{\mathrm{out}}$的经验风险为
\begin{equation}
 \hat R_S(w)=\frac1n\sum_{i=1}^n\ell(w,z_i).
 \label{eq:general-risk}
\end{equation}
\end{definition}

在训练与评价同分布的基线下，我们以$\mathcal W$中的最低期望风险作比较。即使这个基准非零，也可要求输出只比它高出一个预先指定的容差，并让训练失败概率对允许的分布统一受控。

\term{PAC学习}（probably approximately correct learning）将这一要求写成两个容限。$\varepsilon$约束输出风险比比较空间最低风险高出多少，$\delta$约束随机样本与算法随机性导致超出该容限的概率。可学习性要求同一个算法对任意预先选定的$\varepsilon,\delta\in(0,1)$，在足够大的样本量下满足两项要求。

这些比较首先要有定义。下文默认$\inf_{w\in\mathcal W}R_{\mathcal P}(w)$有限，随机输出的风险几乎处处有限，并且$R_{\mathcal P}(A(S,U))$对$(S,U)$可测。无界损失的具体定理还须用可积性或矩条件保证这些性质。满足统计量的定义条件之后，PAC讨论能否学到；程序实现另加计算要求。

\begin{definition}[相对于分布族的PAC可学习性]
\label{def:general-pac-learnability}
若存在算法$A$，对任意$\varepsilon,\delta\in(0,1)$存在整数$n_0(\varepsilon,\delta)$，使任意$\mathcal P\in\mathfrak P$和$n\geq n_0(\varepsilon,\delta)$均满足
\[
 \mathbb P_{S\sim\mathcal P^n,U}\!\left\{
 R_{\mathcal P}(A(S,U))\leq
 \inf_{w\in\mathcal W}R_{\mathcal P}(w)+\varepsilon
 \right\}\geq1-\delta,
\]
则称该学习问题相对于$\mathfrak P$是PAC可学习的。概率包含训练样本和算法内部随机数；确定性算法省略$U$。只有当$\mathfrak P$包含所讨论空间上的全部相关分布时，才称这一保证为分布无关。
\end{definition}

定义~\ref{def:general-pac-learnability}尚未约束程序的执行。若训练在给定编码下停机，且输出的预测程序对每个允许输入停机，得到\term{可计算学习}（computable learning）；再要求样本量、训练时间与预测时间有约定参数的多项式上界，得到\term{高效学习}（efficient learning）。第~\ref{sec:pac-computation-requirements}节会在二分类中落实这两种附加要求。

\section{PAC学习}
\label{sec:pac-learning}

例如，在可实现二分类中，要求“以至少99\%的概率得到未来错误率不超过5\%的分类器”，就是取$\delta=0.01$、$\varepsilon=0.05$。5\%约束固定训练结果在分布上的错误概率，99\%约束随机抽样并训练的成功机会。它既不要求分对每个训练样本，也不要求每批未来邮件都出现至多5\%的实际错误。

有限样本会遗漏输入，所以通常只能允许小概率失败，再用更多样本使失败概率低于指定的$\delta$。成功时的风险应接近零还是接近类内最低值，则由是否存在完美分类器决定。两种情况共享定义~\ref{def:general-pac-learnability}的概率量词；取比较空间$\mathcal H$与零一损失后，可实现与不可知设定分别规定允许的分布及风险基准。

\subsection{可实现情形}

先假设$\mathcal H$中确有一个在总体上零错误的假设。这个可实现条件暂时排除了标签噪声及表示限制带来的不可避免风险，让我们集中分析：未知目标已经在空间内，有限观测是否足以选出低风险假设。

\begin{definition}[可实现性]
\label{def:realizability}
如果某个假设能够以概率1正确预测数据来源中的标签，就称数据分布对该假设空间
是\term{可实现的}（realizable）。形式上，分布$\mathcal P$对$\mathcal H$可实现，
是指存在$h^\star\in\mathcal H$满足
\[
 R_{\mathcal P}(h^\star)=0.
\]
\end{definition}

前面的教学分布由阈值0.6确定真实标签，而$h_{0.6}$属于阈值空间，故该分布可实现。若同一分数可以伴随两种标签，分数本身就未必足以作出零错误预测。这里要求的是总体中的零风险假设存在；一张训练表碰巧没有冲突，只检验了有限观测，不能据此确认可实现性。

存在零风险假设，并没有告诉算法应选哪个。我们要求预先确定的同一算法处理所有可实现分布，充分样本量只依赖$\varepsilon$、$\delta$与$\mathcal H$。若样本量还随未知的具体分布临时确定，就没有得到这种统一保证。

将定义~\ref{def:general-pac-learnability}的允许族限定为对$\mathcal H$可实现的全部分布，记样本量阈值为$m_{\mathcal H}$。此时类内最低风险为零，成功事件化为
\begin{equation}
 \mathbb P_{S\sim\mathcal P^n}
 \left(R_{\mathcal P}(A(S))\leq\varepsilon\right)\geq1-\delta.
 \label{eq:pac-guarantee}
\end{equation}
若同一个算法对所有这些分布达到该要求，就称$\mathcal H$可实现PAC可学习。$m_{\mathcal H}(\varepsilon,\delta)$可以依赖两项容限及假设空间，不能随具体未知分布改变。

在所有算法中，使该保证从某个样本量起对全部更大样本量成立的最小阈值，称为\term{样本复杂度}（sample complexity）。找到一个满足条件的$m_{\mathcal H}(\varepsilon,\delta)$，只证明了它是充分上界，尚未证明它最小。算法随机化时，成功概率也包含内部随机数。

式~\eqref{eq:pac-guarantee}的内层先评价一次训练：$A(S)$是所得分类器，$R_{\mathcal P}(A(S))\leq\varepsilon$就是成功事件，在零一损失下表示总体错误概率不超过$\varepsilon$。外层再对训练样本的随机抽取计数，要求这个事件发生的概率至少为$1-\delta$。

算法与样本量须在未知分布之前确定。若知道分布后才为它挑算法或另定样本量，所得结论只覆盖各自特定来源。分布无关学习要求同一个算法与同一个样本上界覆盖全部允许分布。PAC的计算形式化由Valiant提出；本章采用统计学习教材中的分布无关表述\scite{valiant1984}。

两个容限也约束不同的改进：减小$\varepsilon$使成功结果的风险更低，减小$\delta$使训练失败更少。它们需要增加多少样本，须从假设空间与算法推导。第~\ref{chap:binary-learnability}章将用有限空间给出这种计算。

上面的保证能指定任意小精度，称为\term{强学习}（strong learning）。若只要求成功时比等概率随机猜测好一个固定幅度，则得到\term{弱学习}（weak learning）。在可实现二分类中，要求存在$\gamma\in(0,1/2)$、同一算法$A$和有限$m_{\mathcal H}^{\mathrm{weak}}(\delta)$，对任意$\delta\in(0,1)$、全部可实现$\mathcal P$及$n\geq m_{\mathcal H}^{\mathrm{weak}}(\delta)$满足
\[
 \mathbb P_{S\sim\mathcal P^n,U}\!\left(
 R_{\mathcal P}(A(S,U))\leq\frac12-\gamma
 \right)\geq1-\delta,
\]
这时称$\mathcal H$弱PAC可学习。$\gamma$不随分布、样本或$\delta$变化，$U$独立于样本并表示内部随机性，确定性算法省略它。取$\gamma=0.1$时，成功结果至多40\%错误，比不看输入的公平猜测低10个百分点。第~\ref{sec:boosting}节还将区分目标与输出空间：重加权改变各轮分布，同一学习器与优势必须适用于仍由目标空间实现的每轮分布，并分配总失败预算。一次略优于猜测的训练结果尚未满足这一接口。

\subsection{不可知情形}

如果假设空间没有零风险成员，可实现保证就失去了前提。在应用中，这一前提常无法确认，有时已知不成立。此时应先判断输入和假设空间允许达到什么风险，再评价算法是否达标，不能把未达到零风险直接算作训练失败。

零风险不可达可能来自输入信息不足，也可能来自假设空间遗漏了决策规律。前者即使知道完整分布仍留下错误，后者使类内最低风险高于所有可测分类器的最低风险。第~\ref{sec:statistical-model-selection}节会用贝叶斯风险与近似误差分别刻画二者。噪声的解释也取决于评价标签：若观测标签本身作为$y$，噪声属于$\mathcal P$；若训练与评价用不同标签，须另加二者关系，原同分布保证不能直接使用。

因此，我们改用类内最低风险作为成功基准，同时不对噪声模型预作限制。\term{不可知设定}（agnostic setting）要求结果接近这个基准，而不承诺零风险目标已经在$\mathcal H$内。“不可知”撤去的是这项先验承诺，训练标签与独立同分布抽样仍然保留。

假设$\mathcal H$的最低期望风险是8\%。要求恰当学习器达到2\%总风险就不可能满足；要求它比类内最低风险至多多2个百分点，则允许总风险达到10\%。学习输出的期望风险与类内最低风险之差称为\term{超额期望风险}（excess risk）。它的参照始终是同一个$\mathcal H$，本例的容限为2个百分点。

我们要比较的是可逼近的最低风险，未必存在恰好最优的成员。因此以$R_{\mathcal P,\mathcal H}^\star$记下确界：
\begin{equation}
 R_{\mathcal P,\mathcal H}^\star=\inf_{h\in\mathcal H}R_{\mathcal P}(h).
 \label{eq:best-in-class-risk}
\end{equation}
例如候选风险依次为$0.08+0.1/k$，可趋近0.08却不必取到0.08。下确界保留这个最佳逼近水平，无需额外假定最优假设存在。

再将定义~\ref{def:general-pac-learnability}中的允许族取为$X\times Y$的全部相关联合分布，记充分样本阈值为$m_{\mathcal H}^{\mathrm{agn}}$。相对于上述下确界，要求
\begin{equation}
 \mathbb P_{S\sim\mathcal P^n}\left(
 R_{\mathcal P}(A(S))\leq R_{\mathcal P,\mathcal H}^\star+\varepsilon
 \right)\geq1-\delta.
 \label{eq:agnostic-pac-guarantee}
\end{equation}
满足该统一保证时，就说$\mathcal H$是\term{不可知PAC可学习的}（agnostic PAC learnable）。容限$\varepsilon$限制的是超额期望风险，而非全部分类风险。

可实现分布也在不可知保证覆盖的范围内。其类内最低风险为零，式~\eqref{eq:agnostic-pac-guarantee}便成为可实现保证。因此不可知学习扩大了需要统一处理的分布，而非任意放宽成功标准。类内最低风险只用于评价结果，算法无须事先知道它；这与可实现学习无需知道零风险假设是谁相同。

\subsection{计算要求}
\label{sec:pac-computation-requirements}

前面的PAC定义把算法写作从样本到分类器的映射，回答有限信息是否允许可靠预测。一个满足保证的映射存在，还不足以说明计算机怎样实现它。

实现需要训练和预测都能完成，也需要核算计算量随规模怎样增长。可计算PAC与高效PAC分别加上停机要求和资源上界，统计风险保证仍须成立。

核算之前先固定编码。自然数输入可写为二进制字符串，二分类标签占一个比特，假设可由计算预测的程序表示，学习器输出该程序的有限描述。任意实数却未必有有限精确编码；连续输入上的任意精度运算需要约定精度或计算模型，不能默认按一次常数时间操作计费。

\begin{definition}[可计算PAC学习]
\label{def:computable-pac-learning}
在固定的有效表示下，如果存在一个可执行的学习程序，能够在每个允许的有限训练样本上停机，输出一个对每个允许输入都能停机并给出标签的预测程序，并且该学习算法满足PAC保证，就称该假设空间是可计算PAC可学习的（computably PAC learnable），简称CPAC可学习。这里的PAC保证可以采用可实现版本或不可知版本，使用时必须说明采用哪一种。
\end{definition}

停机要求同时覆盖训练与使用。指定“返回期望风险最小者”没有提供寻找它的程序，输出一段可能永不返回的代码也没有提供总可用的预测规则。允许很慢但最终完成，已经满足可计算性；多项式时间是另一项要求。输出须在$\mathcal H$中时得到恰当CPAC，允许空间外分类器时得到非恰当CPAC。本章默认恰当输出，引用文献结论时须保持输出约定一致。

训练可执行也不等于能算出充分样本上界。前者保证给定有限样本后能完成训练，后者要求根据$\varepsilon,\delta$返回足够大的样本数。基本CPAC未自动要求后者，附加可计算样本上界的版本及分离结果见Delle Rose等\scite{dellerose2023}。第~\ref{sec:computable-binary-learning}节将以可计算缺失标记说明这项区别。

有限穷举可实现，却可能需要指数时间。设假设由$d$比特索引，有$2^d$个成员且各可多项式时间求值。逐一算经验风险仍要检查指数多个候选。定理~\ref{thm:finite-realizable-pac}给出的可实现充分样本量只按下式增长：
\[
 \frac{d\log2+\log(1/\delta)}{\varepsilon}
\]
样本的线性规模与穷举的指数规模并不冲突，它们分别衡量信息与搜索成本。这个比较只排除了穷举的高效性，未排除利用结构的更快算法。

\begin{definition}[高效PAC学习]
\label{def:efficient-pac-learning}
高效PAC学习（efficient PAC learning）要求在满足PAC期望风险保证的同时，样本量、训练计算量及输出分类器的预测计算量都有多项式上界。为明确规模依赖，考虑由参数$d$描述的一族假设空间$\mathcal H_d$，并约定观测编码长度上界为$b$。本书采用的高效性约定要求存在统一程序和固定多项式$p,q$，使所需样本量及训练时间分别满足
\begin{equation}
 \begin{aligned}
 m_{\mathcal H_d}(\varepsilon,\delta)
 &\leq p\bigl(d,b,1/\varepsilon,\log(1/\delta)\bigr),\\
 T_{\mathrm{train}}
 &\leq q\bigl(n,d,b,1/\varepsilon,\log(1/\delta)\bigr).
 \end{aligned}
 \label{eq:efficient-pac-resources}
\end{equation}
学习程序把$d,\varepsilon,\delta$及样本编码作为输入；若$\varepsilon,\delta$只通过预先给定的
样本量体现，也必须明确这一接口。输出分类器属于预先规定、可统一求值的表示族，它在新输入
上的求值时间应被$d$、输出表示长度和输入编码长度的某个固定多项式所控制，不能仅把任意
程序的短代码长度当作预测高效的依据。若这族假设空间还由目标描述长度等量参数化，也须将
它们明确列入规模变量。
\end{definition}

规模参数$d$由问题约定，不一定是特征维数；编码长度$b$使长数值的读写与比较计入成本。统一多项式还要求次数不随$d$增长。不同文献有时使用$1/\delta$，有时使用$\log(1/\delta)$描述资源，引用时须核对。即使满足多项式上界，大常数或高次数也可能使实际计算难以承受。

阈值的顺序结构提供了比穷举更少的搜索位置。设输入为$1,\ldots,M$中的整数，$h_t(x)=\mathbb I\{x\geq t\}$，$t\in\{1,\ldots,M+1\}$。样本上的预测只在阈值越过一个已见输入时改变。将样本排序并合并重复输入后，我们可扫描切分位置，累计左侧正标签数与右侧负标签数，选择错误最少者。排序与扫描分别需要$O(n\log n)$次比较和$O(n)$次累计，每次整数运算另按$\log M$计位成本。预测只需一次阈值比较。结合有限空间的样本界，这个算法在所述编码下满足高效PAC要求。

固定风险目标、表示与输出范围后，高效PAC蕴含可计算PAC，可计算PAC蕴含统计PAC；逆向均须补充条件。可实现与不可知限制数据分布，恰当与非恰当限制输出，可计算与高效限制执行方式。这些要求可以组合使用，不能以其中一项替代另一项。

\section{学习的限制}
\label{sec:limits-of-learning}

对未知$\mathcal P$，算法不能直接比较期望风险，却能逐条计算训练损失。选择经验风险最小者便成为一个可用的训练原则。前面的阈值0.4与0.6在四封邮件上都零错误，也都可能按这个原则被选中。

\begin{definition}[经验风险最小化]
\label{def:erm-learning-rule}
\term{经验风险最小化}（empirical risk minimization, ERM）是在给定假设空间中，选择
训练样本平均损失最小的假设作为学习结果。用$A_{\mathrm{ERM}}$表示这种学习算法，则
\begin{equation}
 A_{\mathrm{ERM}}(S)\in\argmin_{h\in\mathcal H}\hat R_S(h).
 \label{eq:erm-learning-rule}
\end{equation}
这里$\argmin$表示所有达到最小经验风险的假设组成的集合。若最小值不能取得，可以
要求输出的经验风险与下确界相差不超过一个给定容差。
\end{definition}

$\argmin$是一组经验最优解，公式中的“$\in$”允许算法选其中任意一个。ERM本身尚未指定并列解怎样选择。两条阈值的经验风险相同而未见预测不同，说明要从经验最优推到PAC成功，必须另行控制未观察区域的风险。

扩大样本也未必消除这种差别。若空间允许按每份样本特制记忆规则，新增数据可以只使记忆表变长。下面用连续输入上的有限记忆规则，检验零训练错误能否单独推出低期望风险。

\subsection{记忆训练样本}

仍以$[0,1]$中的邮件分数作输入，这次令它均匀分布，区间概率等于长度，并令所有真实标签都为垃圾邮件。这是人为简化的分布，便于直接比较记忆与推广。

始终预测垃圾邮件的规则在此分布上零错误。另一个规则只保存训练分数：查到记忆表中的分数时预测垃圾邮件，表外分数则预测正常邮件。二者都能分对训练邮件，但表外预测相反。

将记忆表写成有限分数集合$C$，规则为
\begin{equation}
  h_C(x)=
  \begin{cases}
    1, & x\in C,\\
    0, & x\notin C,
  \end{cases}
  \label{eq:finite-memory-classifier}
\end{equation}
再令$h_{\mathrm{all}}(x)=1$，把它与全部有限表规则$h_C$一起放入假设空间。每条规则可测，且$h_{\mathrm{all}}$的总体风险为零，故这个分布可实现。

以$C_S$记录样本$S$出现过的全部不同输入。选择$h_{C_S}$是一个合法的零经验风险选解，但它的两项风险为
\begin{equation}
  \hat R_S(h_{C_S})=0,
  \qquad
  R_{\mathcal P}(h_{C_S})=1.
  \label{eq:memorization-erm-failure}
\end{equation}
每个训练点都在表内，所以第一项为零；连续均匀分布落在任意有限集合的概率为零，所以新输入几乎必在表外，第二项为1。零训练错误与几乎处处预测错误在同一规则上同时成立。

始终选$h_{\mathrm{all}}$的ERM在这个例子中仍能成功，反例只证明ERM原则不能单独约束选解。它要求已见点上的损失最小，却允许在表外任意回答。第~\ref{sec:binary-uniform-convergence}节将用一致收敛连接经验风险与总体风险，使所有经验最优解都接近类内最低期望风险。

\subsection{没有免费的午餐}
\label{sec:no-free-lunch}

记忆规则失败后，一个自然的问题是：换一套选择规则，能否对任意标签规律都保证成功？如果假设空间允许许多输入的标签彼此无约束地变化，未见标签就没有可由已见标签推断的联系。\term{没有免费的午餐}（no-free-lunch）定理把这一信息缺口写成适用于任意有限样本算法的失败下界。

设算法只看到$n$个观测，而空间能实现某$2n$个输入上的全部二元标签组合。我们将固定输入出现的概率，再为算法选择一套标签规律。这套规律在$\mathcal H$中有零风险实现，算法却以至少$1/7$的概率得到期望风险不低于$1/8$的结果。下面的构造解释两个数值从哪里来。

算法不知道完整输入分布及标签表，只观察抽到的$n$个记录。$2n$个输入无法全部出现，未见点的标签又可以独立改变，所以拟合已见数据不能排除在未见点大量出错的可能。这里用总体风险超过容限判定失败，训练错误仍可为零。

先用四个输入观察歧义。令$X=\{1,2,3,4\}$且四者等概率出现，编号仅用来区分对象，不提供类别或相似性信息。假设空间允许全部二分类函数。若训练只出现$(1,0)$与$(2,1)$，编号3与4的标签仍然未见。

完整标签表$(0,1,0,0)$与$(0,1,1,1)$分别记为$f_0,f_1$。它们在训练点上完全相同，都解释已见标签，在编号3、4上却处处相反。因此“与样本一致”只表示在观察到的点上回答正确。图~\ref{fig:learning-no-free-lunch}并列画出这两套隐藏真值。

\begin{figure*}[!htb]
  \centering
  \includegraphics[width=164.600mm]{figures/03-foundations/02-learning-theory/generated/learning-theory-no-free-lunch/figure.png}
  \caption{相同样本与两套隐藏标签。鸟、汽车、狗、船分别表示编号1至4，每个概率为$1/4$；两套来源在已见编号1、2上同为0、1，未见编号3、4则分别全为0或全为1。图示输出$h(3)=0,h(4)=1$在每套来源各错一个点。固定样本与内部随机数后输出相同，任意二元预测在每个未见点的两份损失相加为1，故风险之和至少$1/2$。这里展示固定样本的歧义，下文再证明随机抽样下的固定困难分布。}
 \label{fig:learning-no-free-lunch}
\end{figure*}

预测编号3时，$f_0$要求回答0，$f_1$要求回答1。两套真值交给算法的样本都是$(1,0),(2,1)$，算法无法用这份样本辨认当前是哪一套。它可以按偏好选择一个标签，但该偏好仍可能与真实标签相反。

把所有假设枚举出来仍只能发现$f_0,f_1$都零训练错误。枚举提供了候选规则，没有提供二者在编号3上谁正确的证据。要消除歧义，须观察该标签，或加入已见与未见标签之间的联系。

同样的歧义可放在$2n$个等概率输入上。给定$n$个观测，未见集合$U$至少含$n$点。两张表在已见处相同、在$U$上相反时，同一输出$h$在每个未见点的两份零一损失相加为1，故
\[
 R_{\mathcal P_0}(h)+R_{\mathcal P_1}(h)
 \geq\frac{|U|}{2n}\geq\frac12.
\]
至少一张表使这次输出的风险不低于$1/4$。此处仍在固定样本下对照候选；下文须把候选平均，才能在抽样之前选定一个困难分布。

固定样本下的两套标签表展示了信息歧义，但学习定理要求固定分布在随机抽样时仍然困难。为满足这个量词顺序，我们将对全部标签配置平均，再从中选定一个分布\scite{shalevshwartz2014}。

\begin{theorem}[没有免费的午餐]
\label{thm:no-free-lunch-classification}
设$n\geq1$，固定输入空间$X$中有一个含$2n$个不同点的集合$C$。假设无论怎样为
$C$中的每个点分别指定标签0或1，都能从$\mathcal H$中找到一个分类器，使它在$C$上的
预测与这些指定标签完全一致；不同的标签组合可以由$\mathcal H$中的不同分类器实现。
则对任意使用$n$个独立同分布训练观测的二分类学习算法$A$，都存在一个对$\mathcal H$
可实现的分布$\mathcal P$及$f^\star\in\mathcal H$，使$R_{\mathcal P}(f^\star)=0$，但
\begin{equation}
 \mathbb P\left(R_{\mathcal P}(A(S))\geq\frac18\right)\geq\frac17.
 \label{eq:no-free-lunch-probability}
\end{equation}
这里概率对$S\sim\mathcal P^n$及算法内部随机性计算。算法即使允许输出$\mathcal H$以外的分类器，结论也成立。
\end{theorem}

\begin{proof}
固定输入集合$C$并让输入均匀出现，候选分布之间只改变标签表。对任意$g:C\to\{0,1\}$，前提给出$f_g\in\mathcal H$与它在$C$上一致。按先抽输入、再取标签$g(x)$的方式定义$\mathcal P_g$，这个分布就由$f_g$实现，风险为零。

暂时给$C$中每个输入独立分配一个公平的二元标签，用来对候选标签表平均。每个点只分配一次，重复抽到它时标签保持不变。因此任意固定标签表仍对应无噪声的可实现分布；随机标签表是证明工具，不是训练过程中不断变化的真值。

固定训练输入、已见标签与算法随机数后，算法输出也已确定。$n$个观测至多涉及$n$个不同输入，故$C$中至少$n$个输入未见。每个未见标签仍是独立公平的二元选择，无论输出在该点回答0还是1，条件错误概率都是$1/2$。

每个未见输入占概率$1/(2n)$。风险非负，只累加至少$n$个未见点已给出随机标签表下的下界：
\begin{equation}
 n\times\frac1{2n}\times\frac12=\frac14.
 \label{eq:no-free-lunch-expected-risk}
\end{equation}
三个因子依次为未见点数下界、单点概率与条件错误概率。再平均训练输入、已见标签和内部随机数仍有$\mathbb E_g\mathbb E_{S,U}R_{\mathcal P_g}(A(S,U))\geq1/4$，其中$U$表示内部随机数。

候选表只有有限张，平均至少为$1/4$就必有固定$g^\star$满足
\[
 \mathbb E_{S,U}R_{\mathcal P_{g^\star}}(A(S,U))\geq\frac14.
\]
若每张表的平均风险都小于$1/4$，总平均也会小于它，产生矛盾。现在固定$\mathcal P=\mathcal P_{g^\star}$及$f^\star=f_{g^\star}$，后续概率只包含从这个固定来源抽样与算法随机性。

最后把平均风险下界转成失败概率。令$Z=R_{\mathcal P}(A(S,U))\in[0,1]$、$p=\mathbb P(Z\geq1/8)$。失败事件外$Z<1/8$，事件内$Z\leq1$，从而
\[
 \frac14\leq\mathbb E Z\leq(1-p)\frac18+p
 =\frac18+\frac78p.
\]
移项即得$p\geq1/7$。\end{proof}

困难分布是在算法$A$与样本量$n$固定后选出的。此后一切重复训练都从这个同一分布抽样，输入频率与标签表不再改变；以至少$1/7$的概率，输出风险不低于$1/8$。证明并未在看到某次样本后再改标签，因而给出的是固定来源上的随机训练失败，而非一份事后挑出的不利样本。

分布可以随$A,n$变化，因为定理对每个算法及每个给定样本量分别构造反例。它没有要求同一个分布击败所有算法，也没有声称某算法在每个分布上都失败。算法仍可能处理许多有结构的来源；被排除的是在这些标签完全自由的允许分布上统一成功。

将这一有限样本结论用于固定的无限空间$X=\mathbb N$，令$\mathcal H$包含全部二分类函数。假设它PAC可学习，取$\varepsilon=1/16$、$\delta=1/8$，便应有有限的统一充分样本量$n\geq1$。从同一个$X$选出$2n$个点，定理给出的可实现分布使风险至少为$1/8$的概率不低于$1/7$，同时超过所选误差与失败容限，形成矛盾。不能统一学习的是这一不受限制的空间，不是每个具体分布。

可实现性保证完整真值表$f^\star$在$\mathcal H$内，却未把这张表交给学习器。学习器只见$n$个带标签记录。某编号若未出现，将其标签改成另一类而保持全部已见标签，就能得到同样解释样本的另一张表。零风险解存在与样本足以辨认其未见预测，是两项不同条件。

这也说明失败并非搜索时间不足。即使允许无限计算，算法依然看到同一份样本及彼此冲突的候选真值，无法靠搜索判断哪张表才是真实来源。第~\ref{sec:pac-computation-requirements}节的不可计算或时间过长限制执行；这里的限制来自缺少标签信息。

PAC无需复原整张真值表，只要求总体风险低于$\varepsilon$。上面的下界却已表明：任意有限统一样本量，都存在允许分布使该风险要求以过大的概率失败。因而即使采用这个较弱的预测目标，也不能为不受限制的空间找到统一的有限样本阈值。

正面保证需要让一个标签约束其他输入的预测。例如，若阈值规则分对标签为1、分数为0.5的邮件，必有$t\leq0.5$，所有更高分数也就被预测为1。观察一个标签排除了许多阈值，并限制了未见预测。若每个编号的标签都可自由选择，编号1的标签便不能提供编号2的任何信息。

限制假设空间或分布，正是在加入这种可以推广的联系。限制也有代价：若真实规律不满足它，所需假设会被排除。下一节在同一个风险口径下分析这项表达代价与有限样本的选择代价。

\section{风险来源与模型选择}
\label{sec:statistical-model-selection}

若垃圾邮件的分数落在两个分离区间，中间却是正常邮件，单阈值便不能同时圈出两段正类而排除中段负类。阈值限制能减少自由预测，却也留下训练数据无法消除的表达误差。要选择空间，须将这种误差与输入信息本身的不足分开。

输入信息不足可以在一个分数处看清。设分数0.5的邮件中70\%为垃圾、30\%为正常，分类器只看到$x=0.5$，所以只能给出同一个预测。预测垃圾会错30\%，预测正常会错70\%，任意分类函数也无法低于30\%的条件错误率。只要表示仍为该分数，更多样本或另一算法就不能区分这些邮件。

空间限制则是另一种情形。设三个输入分数为0.2、0.5、0.8，标签分别确定为1、0、1。输入足以区分它们，也存在全部分对的函数，但单阈值无法产生1、0、1的变化，每条阈值至少错一类。此时缺少的不是区分输入的信息，而是$\mathcal H$中能利用这种信息的规则。

因此，先求当前输入下所有分类器的最低风险，再看选定空间使这个最低值增加多少。第一个基准由贝叶斯风险给出，第二个差距度量空间的表达限制。

\term{贝叶斯风险}（Bayes risk）是在当前输入信息下，全部可测分类器的最低错误概率。在条件风险有限、最小值达到且选择可测时，按给定结果分布选择条件期望损失最小的行动，就得到Bayes决策；这里采用二分类零一损失（可参考\BookSectionRef{sec:decision-bayes-rules}{第一卷“策略的价值与优化”章的“单步价值下的行动选择”一节}）。令$\mathcal H_{\mathrm{all}}$为全部可测二分类器，便有共同基准
\begin{equation}
 R_{\mathcal P}^\star=\inf_{h\in\mathcal H_{\mathrm{all}}}R_{\mathcal P}(h).
 \label{eq:bayes-risk}
\end{equation}
这个基准不限制分类器属于所选$\mathcal H$，表示知道数据规律时可达到的水平，没有假定学习器已经知道它。

求共同基准可以先在一个输入$x$处比较两种标签。设$X$为标准Borel空间；更一般地，要求$\mathbb E[\mathbb I\{y=1\}\mid x]$有可测版本。该版本按$\mathcal P_X$几乎处处确定，选一个并记为
\[
 \eta(x)=\mathbb P(y=1\mid x),
\]
它记录输入$x$处标签为1的条件概率，标签为0的概率相应为$1-\eta(x)$。

在$x$处预测1的错误概率为$1-\eta(x)$，预测0的错误概率为$\eta(x)$。选择较小错误率，就应在$\eta(x)\geq1/2$时预测1，否则预测0；相等时两者均最优。逐点这样选择得到\term{贝叶斯分类器}（Bayes classifier），用所选可测版本写为
\[
 h^\star(x)=\mathbb I\{\eta(x)\geq1/2\}.
\]
每个输入处的最小条件错误率因而为
\[
 \min\{\eta(x),1-\eta(x)\}.
\]

在分数0.5处，$\eta(0.5)=0.7$意味着两种预测的条件错误率分别为0.3与0.7。贝叶斯分类器选择垃圾邮件，仍错其余30\%的正常邮件。这个例子将“选多数标签”的规则与只观察分数时的不可避免错误联系起来。

条件错误率要按输入频率加权才能得到总体风险。例如分数0.5占全部邮件的$q$，0.3的条件错误贡献$0.3q$。由联合分布忽略标签得到输入\term{边缘分布}（marginal distribution）$\mathcal P_X$，用它平均上述逐点最小值，得到
\[
 R_{\mathcal P}^\star=
 \mathbb E_{x\sim\mathcal P_X}[\min\{\eta(x),1-\eta(x)\}].
\]
若两类标签都有正条件概率的输入集合具有正$\mathcal P_X$测度，这个平均便为正。即使已知分布、采用贝叶斯规则，当前输入仍会留下错误。

\subsection{两类误差}

相对于当前输入的贝叶斯风险，学习结果多出的风险可分成两段：限制候选空间使最低风险升高，算法在这个空间中的实际选择再使风险偏离类内最低值。

第一段称为\term{近似误差}（approximation error），等于类内最低期望风险减去贝叶斯风险。它由候选空间决定；固定空间后，即使拥有无限数据并能找到类内最优者，也不能因此消除该差距。

第二段通常称为\term{估计误差}（estimation error），等于输出的期望风险减去类内最低期望风险。对任意算法，它更准确地表示\emph{类内超额风险}，还可能包含优化未完成、正则化改变目标或不一致选解的影响。对ERM、近似ERM或另有一致性保证的算法，才能将相应上界主要解释为有限样本的统计估计代价。

对恰当输出$\hat h=A(S)\in\mathcal H$，在贝叶斯风险与输出风险之间加入类内下确界，便得
\begin{align}
 R_{\mathcal P}(\hat h)-R_{\mathcal P}^\star
 ={}&\left(\inf_{h\in\mathcal H}R_{\mathcal P}(h)-R_{\mathcal P}^\star\right)
 \label{eq:approximation-estimation-decomposition}\\
 &+\left(R_{\mathcal P}(\hat h)-\inf_{h\in\mathcal H}R_{\mathcal P}(h)\right).
 \notag
\end{align}
第一项是近似误差，第二项是估计误差。$\mathcal H\subseteq\mathcal H_{\mathrm{all}}$使第一项非负，$\hat h\in\mathcal H$使第二项非负。允许空间外输出时，它可以优于$\mathcal H$的下确界，第二项便可能为负，原非负分解的解释也须相应改变。

取一个教学数值例子：贝叶斯风险5\%，阈值空间最低风险12\%，所学阈值风险15\%。表达限制贡献7个百分点，类内选择贡献3个百分点，总共高于贝叶斯风险10个百分点。具有一致性保证的算法可用更多数据缩小第二段的统计部分，始终返回坏假设的算法却不会自动改善。第一段的7个百分点还须通过改变输入表示或扩大假设空间处理。

图~\ref{fig:approximation-estimation-error}沿风险增大的方向排列这三个水平。贝叶斯风险到类内最低风险的距离是近似误差，类内最低风险到输出风险的距离是估计误差。后一个距离的来源仍取决于算法；只有结合ERM、近似ERM或其他一致性保证，才能用有限样本代价解释其上界。

\begin{figure}[htbp]
  \centering
  \includegraphics[width=112.000mm]{figures/03-foundations/02-learning-theory/matplotlib/learning-theory-error-decomposition/figure.pdf}
  \caption{风险的两段差距。教学设定中，贝叶斯风险为5\%，阈值空间最低风险为12\%，输出风险为15\%。前两者相差7个百分点，为实线表示的近似误差；后两者相差3个百分点，为虚线表示的估计误差。}
  \label{fig:approximation-estimation-error}
\end{figure}

扩大候选空间提供了更多逼近目标的机会，也提供了更多偶然拟合样本的机会。\term{偏差与复杂度的权衡}（bias--complexity trade-off）指的就是二者之间的选择：“偏差”在这里表示空间的表达限制，“复杂度”表示有限样本中可靠选解的统计代价。

若正类由两个分离区间组成，单阈值的近似误差不会因数据无限增加而消失。允许许多小区间以后，空间可以表达这些正类，也可以紧密包围训练点；未见间隙仍可能被误判，类内选择的统计误差便更难控制。式~\eqref{eq:approximation-estimation-decomposition}的两项随空间一起变化，最终风险未必下降。结构风险最小化将据此比较拟合收益是否足以补偿新增的不确定性。

算法对训练目标的求解差距还须单列。若样本上最好候选错2封、返回者错5封，多出的3次错误来自经验目标未达到最优。\term{经验优化误差}（empirical optimization error）以经验风险之差计量：
\begin{equation}
  \Delta_{\mathrm{opt}}(S)
  =\hat R_S(\hat h)-\inf_{h\in\mathcal H}\hat R_S(h)
  \geq0,
  \label{eq:empirical-optimization-error}
\end{equation}
它比较样本目标，与式~\eqref{eq:approximation-estimation-decomposition}的期望风险差不同，不能把二项恒等式直接变成三个非负期望风险差。须先用泛化控制连接经验与总体，再让$\Delta_{\mathrm{opt}}(S)$进入输出风险保证。

\subsection{用验证集选择模型}
\label{sec:validation-model-selection}

若嵌套的空间逐渐扩大且每次精确最小化经验风险，训练误差只能持平或下降。这种下降无法单独判断未来误差，可以用不参与参数拟合的\term{验证集}比较复杂度或其他超参数。第~\ref{sec:experience-data-boundary}节说明训练、验证、测试各自的操作边界。候选越多，偶然低估的验证误差越容易被选中，最终评价因而需要覆盖“调参并选择模型”的整个流程（可参考\BookSectionRef{sec:statistics-bias-efficiency}{第一卷“随机变量”章的“估计误差与效率”一节}）。

图~\ref{fig:learning-model-selection}以教学曲线示范：训练误差继续下降时，验证误差可以回升，所选复杂度由最低验证误差确定。实际曲线未必只有一个谷底。验证数据已用于选择，故锁定模型及超参数后还须以独立\term{测试集}评价最终结果。若继续按测试误差改模型，测试数据也参与了选择，就失去这次独立评价的含义。

\begin{figure}[htbp]
  \centering
  \includegraphics[width=112.000mm]{figures/03-foundations/02-learning-theory/matplotlib/learning-theory-model-selection/figure.pdf}
  \caption{教学示意中的验证选择与独立测试。实线圆点为训练误差，虚线方块为验证误差，圆环的最低验证误差确定$\hat k$。三角形仅为锁定模型后的一次测试评价，未参与选择，故不连成测试曲线。所有数值均为构造示意。}
 \label{fig:learning-model-selection}
\end{figure}

\subsection{结构风险最小化}

把候选空间写成$\mathcal H_1\subseteq\mathcal H_2\subseteq\cdots$，例如$\mathcal H_k$允许至多$k$个区间的并作为正类区域。纯粹比较最小经验风险会偏向更丰富的空间，因为它至少保留前一空间的所有候选；只选最简单的空间，又可能留下较大的近似误差。我们因此需要一条同时考虑表达限制与有限样本选择代价的规则。

\term{结构风险最小化}（structural risk minimization, SRM）以带惩罚的经验目标选择空间，因为近似与估计误差依赖未知分布，不能直接最小化。先固定总失败预算$\delta$，分配与样本无关的$\delta_k>0$、$\sum_{k\geq1}\delta_k\leq\delta$。以$\operatorname{pen}(k,n,\delta_k)$记录第$k$个空间的复杂度惩罚，选择
\begin{equation}
  (\hat k,\hat h)
  \in\argmin_{k\geq1,\,h\in\mathcal H_k}
  \left\{\hat R_S(h)+\operatorname{pen}(k,n,\delta_k)\right\}.
  \label{eq:structural-risk-minimization}
\end{equation}
经验项衡量训练拟合，惩罚控制由样本推断总体的不确定性，通常随空间复杂度增加。它是估计误差的控制量，并非该误差本身。惩罚偏小可能选中偶然拟合的复杂规则，偏大可能留下简单空间的表达偏差。最小值不可达时，可采用带惩罚目标的近似最优解。

同时比较空间，需要同时有效的偏差保证。每个空间单独失败概率不超过$\delta$，并不保证看过数据再选空间的总失败概率仍不超过$\delta$。预先分配$\delta_k$并令$\sum_k\delta_k\leq\delta$后，并集界就控制了至少一项估计失准的概率。因此以至少$1-\delta$的概率，所有空间同时满足各自的偏差要求。

\begin{theorem}[结构风险最小化的期望风险保证]
\label{thm:srm-oracle-bound}
设$\mathcal H_1\subseteq\mathcal H_2\subseteq\cdots$是预先给定的非空假设空间，损失为零一损失，训练样本$S\sim\mathcal P^n$。固定$\delta\in(0,1)$，选择与样本无关的$\delta_k>0$，满足$\sum_{k\geq1}\delta_k\leq\delta$。假设非负惩罚$p_k=\operatorname{pen}(k,n,\delta_k)$使每个假设空间都满足
\begin{equation}
 \mathbb P_{S\sim\mathcal P^n}\left(
 \sup_{h\in\mathcal H_k}
 |R_{\mathcal P}(h)-\hat R_S(h)|>p_k
 \right)\leq\delta_k.
 \label{eq:srm-classwise-confidence}
\end{equation}
这里假定相关事件可测，且惩罚在观察样本前已经确定。

若学习算法输出$\hat h\in\mathcal H_{\hat k}$，并将带惩罚的经验目标求解到给定容差$\rho\geq0$，即
\begin{equation}
 \hat R_S(\hat h)+p_{\hat k}
 \leq\inf_{k\geq1,\,h\in\mathcal H_k}
       \{\hat R_S(h)+p_k\}+\rho,
 \label{eq:srm-approximate-minimizer}
\end{equation}
则以至少$1-\delta$的概率，
\begin{equation}
 R_{\mathcal P}(\hat h)
 \leq\inf_{k\geq1}\left\{
       \inf_{h\in\mathcal H_k}R_{\mathcal P}(h)+2p_k
       \right\}+\rho.
 \label{eq:srm-oracle}
\end{equation}
精确最小化带惩罚的经验风险时取$\rho=0$。若每个假设空间的这些偏差概率界对所有
允许分布都使用同一组惩罚成立，则结论也对这些分布统一成立。
\end{theorem}

\begin{proof}
先把各空间的偏差保证合在一个共同事件中，再在事件内比较输出与任意参照假设。

令$E_k$表示$\mathcal H_k$中全部$h$满足$|R_{\mathcal P}(h)-\hat R_S(h)|\leq p_k$的事件，令$E=\bigcap_{k\geq1}E_k$。并集界把至少一个空间失准的概率控制为
\[
 \mathbb P(E^{\mathrm c})
 =\mathbb P\left(\bigcup_{k\geq1}E_k^{\mathrm c}\right)
 \leq\sum_{k\geq1}\mathbb P(E_k^{\mathrm c})
 \leq\sum_{k\geq1}\delta_k\leq\delta.
\]
事件都由同一份样本决定，无需相互独立。已有$\delta_k$预算便给出$\mathbb P(E)\geq1-\delta$。

固定$S\in E$后，偏差界也覆盖数据选出的空间$\mathcal H_{\hat k}$及输出$\hat h$，因此
\[
 R_{\mathcal P}(\hat h)\leq\hat R_S(\hat h)+p_{\hat k}.
\]
以任意$k$及$h\in\mathcal H_k$作参照，带惩罚目标的近似最优性给出
\[
 \hat R_S(\hat h)+p_{\hat k}
 \leq\hat R_S(h)+p_k+\rho.
\]
对同一参照，$E$又保证$\hat R_S(h)\leq R_{\mathcal P}(h)+p_k$。合并这三次比较，得到
\[
 R_{\mathcal P}(\hat h)
 \leq R_{\mathcal P}(h)+2p_k+\rho.
\]
第一个$p_k$来自参照目标的显式惩罚，第二个来自参照经验风险与总体风险的偏差。它们共同上界控制空间选择的统计代价，并非估计误差的定义值；输出空间的$p_{\hat k}$已经在带惩罚比较中消去。

任意$h\in\mathcal H_k$与任意$k$都满足上式，依次取两层下确界即得式~\eqref{eq:srm-oracle}，无需最优假设或空间实际存在。由于全程在$E$内推导，而$\mathbb P(E)\geq1-\delta$，得到所需高概率结论。
\end{proof}

算法只能计算带惩罚的经验目标，定理却允许它与各空间的最低期望风险加统计代价比较。这类结果称为\term{预言机不等式}（oracle inequality）：参照选择好像知道各空间的总体风险，实际算法仍只使用样本。这里引入预言机的是评价基准，没有增加训练可见信息。

可用望远镜求和分配预算：取$\delta_k=\delta/[k(k+1)]$，因$1/[k(k+1)]=1/k-1/(k+1)$而总和为$\delta$。若各$\mathcal H_k$有限，第~\ref{sec:binary-uniform-convergence}节的一致收敛界给出
\[
 p_k=\sqrt{\frac{\log(2|\mathcal H_k|/\delta_k)}{2n}}
 =\sqrt{\frac{\log(2|\mathcal H_k|)+\log(k(k+1)/\delta)}{2n}}.
\]
$|\mathcal H_k|$计入类内选择代价，$\delta_k$还计入跨空间比较代价。无限空间须另用容量工具建立式~\eqref{eq:srm-classwise-confidence}，不可把无限基数直接代入有限空间的计数公式。

从式~\eqref{eq:srm-oracle}两侧减去贝叶斯风险，每个参照项就成为该空间的近似误差加$2p_k$，另有求解容差$\rho$。$2p_k$上界控制选择的统计代价，不等于估计误差本身。小空间可能留下大偏差，大空间通常需要大统计项；定理用其中最小的总和控制输出风险。它既未承诺零风险，也未提供带惩罚目标的高效求解程序。

\section{研究内容}
\label{sec:learning-theory-research-content}

PAC、不可能性与模型选择已经说明统计保证要同时指定什么信息和什么风险。后续理论又改变标签、损失、算法与环境，所需条件随之扩展。下面沿分析对象的变化梳理研究历史，再把正在研究的问题放入共同地图，以便识别各条路线分别补上了哪项条件。

\subsection{发展历史}

先看固定规则与数据选出的规则有何区别。统计推断研究有限样本、估计误差与重复抽样；模式识别和人工智能还要求算法从许多候选中选择预测规则。一条预先固定规则的经验频率收敛，并不自动保证看过数据后选出的规则同样可靠。学习理论因而需要分析这种选择新增的偏差，而不只研究一个固定估计量。

Vapnik与Chervonenkis将有限输入上的划分数与经验频率一致逼近总体概率联系起来\scite{vapnik1971}。计数把候选函数类的丰富程度变成可分析条件，概率界则控制同一份样本对整个类的评价偏差。二者结合形成了VC维与一致收敛的分类框架，使风险保证能够覆盖数据选出的分类器。

Valiant在1984年以PAC学习明确算法、样本量、误差与失败概率的量词关系，同时要求关注多项式样本量及运行时间\scite{valiant1984}。评价于是从极限收敛推进到允许问题族上的统一有限样本保证。围绕这一标准，有限空间、VC维、ERM与样本复杂度相互联系，统计学习器存在、程序可计算和程序高效也逐步成为分别研究的要求。

这些参数给“学会”提供了可比较的口径：精度限定风险偏离，置信参数限定失败机会，样本复杂度计算满足二者所需的观测数。量词顺序仍决定标准强弱。每个分布各配一个算法，不等于同一算法覆盖全部分布。本章前面的PAC定义保留了后一种统一要求。

在这套口径下，少量样本足以辨认低风险规则，并不自动提供找到规则的高效程序；迅速降低训练损失，也不自动降低期望风险。计算学习理论因此分析表示、查询、归约与运行时间，分别核对信息是否足够及程序能否利用信息。非凸优化、大规模训练和组合输出问题仍会遇到这两项不同的限制。

当评价涉及实值预测与替代损失时，二元划分计数不足以记录函数值的变化。Rademacher复杂度用随机符号检验函数类在样本上的拟合能力\scite{bartlettmendelson2002}，伪维、胖打散维与覆盖数则从不同角度推广容量控制。样本压缩改问能否以少量代表观测和有限旁信息重构预测规则\scite{moran2016compression}。多分类进一步显示，一致收敛、一般可学习性及具体ERM样本复杂度可受不同维度控制，标签结构与输出范围必须计入结论\scite{daniely2015multiclass,brukhim2022multiclass}。

这些扩展保留了原来“同时评价候选”的问题，增加了评价可利用的信息。VC维取最坏输入集合上的二元划分丰富程度，实值函数与间隔分析还需要样本几何、预测分值或损失变化。覆盖、随机复杂度和间隔界分别在这些对象上建立控制，使研究能够在分布无关保证与更贴近数据的保证之间选择。

容量进入模型选择时，就形成经验拟合与复杂度惩罚的比较。表达不足留下近似误差，有限样本中的偶然拟合增加选择代价，经验目标未充分求解又引入优化误差。将三者分别计算，才能看清扩大模型、增加数据或改进求解各自影响哪个环节，也能连接理论界与模型、训练算法的设计。

压缩与多分类还要求核对容量工具到底提供充分条件还是完整刻画。少量代表样本可通过重构限制可传递信息，因而给出泛化保证；不同标签结构下，刻画一致收敛的维度却未必刻画全部学习器。二分类中的等价关系因此不能未经证明推广到任意输出空间\scite{moran2016compression,daniely2015multiclass,brukhim2022multiclass}。

\Needspace{18\baselineskip}
若整个候选类很大，单看它的容量仍无法区分具体算法如何选解。算法稳定性考察删除或替换一个观测后输出的变化，Bousquet与Elisseeff建立了它与泛化的联系；后续研究将这一思路用于一般学习和随机梯度下降的扰动传播\scite{bousquetelisseeff2002,shalevshwartz2010learnability,hardt2016sgd}。信息论则分析输出对训练数据的依赖，以互信息、条件互信息及换测度工具控制期望泛化差\scite{xu2017information}。两条路线都把算法实际使用样本的方式纳入保证。

固定同一个函数空间，正则化ERM、随机梯度下降和其他选解规则可能对样本替换作出不同反应。稳定性把这种反应保留在分析中，限制单个观测的偶然性怎样影响预测。不过，始终返回同一决策也可以稳定，风险却可能很高。稳定性必须结合经验优化或风险基准，才能支持学习成功的结论。

信息论以互信息、相对熵等量记录这种数据依赖。普通互信息衡量输出与整份样本的平均联系，条件互信息在预先抽取的双样本上记录输出透露了哪些选择；压缩长度与有限输出范围还能提供信息量上界。这类分析通常先控制期望泛化差。转成PAC的高概率保证时，还须另加集中、尾界或期望到概率的转换。

显式正则化、样本顺序、随机初始化与提前停止，都能改变同一函数空间中的选解。容量界若只记录空间，就无法呈现这些差别；稳定性与信息量则分别记录算法的样本敏感性及输出保留的数据依赖。研究复杂度的对象由此从候选函数类扩展到算法的行为。

扩展分析对象也要限制结论。很少利用数据的算法可能既稳定又低互信息，却没有学好任务。控制数据依赖只能回答泛化偏离，还需经验拟合、近似能力与比较基准才能判断总风险。深度学习理论同样须分开分析表示形成、优化成功与最终风险，不能让一个依赖量承担全部解释。

这种分析需要保留相应训练条件。容量通常在训练前给出，稳定性与信息依赖却还受随机化方式、运行步数及输出规则影响。初始化、样本顺序、轨迹与停止规则因此成为区分训练过程的材料。后续网络理论据此研究优化轨迹、参数分布和表示变化，使复杂度随模型空间与训练过程共同确定。

\Needspace{20\baselineskip}
现代网络甚至可以拟合随机标签，因此表达能力和参数数量不能单独解释为何某次训练能泛化\scite{zhang2021rethinking}。分析转向所选解的权重范数与间隔、初始化附近的神经切线核、梯度下降的隐式偏差、宽网络平均场及参数分布的PAC--Bayes界\scite{jacot2018ntk,soudry2018implicit,maurer2004pacbayes}。函数、轨迹与参数分布各保留不同信息，分别服务于表示、优化、选解与泛化这些相关而不同的问题。

\Needspace{16\baselineskip}
过参数化模型也可能在近零训练误差后，随规模增加出现测试风险再下降。线性插值中的良性过拟合结果表明，协方差谱、噪声及最小范数选解满足特定条件时，可以控制这类风险\scite{belkin2019double,bartlett2020benign}。因此“参数必须少于样本”不能作为普遍标准，原始参数数目还未记录数据几何、优化偏好与有效自由度。需要判断的是哪类插值解能够控制噪声、保留推广所需结构。

各条网络路线的对象与条件仍不同。神经切线核在初始化邻域中以初始梯度所决定的核近似网络；隐式偏差考察无显式正则化时优化偏向哪些解；平均场研究宽度增加时参数分布的演化；PAC--Bayes比较数据相关参数分布与预定参照分布。它们可在特定模型上相互印证，结论却须各按条件使用。显著特征学习离开初始化邻域时，局部线性化结果尤其不能直接沿用。

预训练又增加了下游学习的前置经验。语料、架构、目标与优化共同决定初始表示，下游样本继续调整，推理还可能由上下文示例与额外计算改变。下游风险因而同时反映预训练覆盖、表示充分性、适应方法与评价协议。只数微调参数不足以分辨这些贡献，分析须分别追踪预训练、迁移和推理阶段提供了什么信息。

\Needspace{8\baselineskip}
训练与评价来源改变后，领域自适应和领域泛化要求重新连接两个环境。分布差异、覆盖、共享标签机制与共同错误等条件提供这种联系；缺少联系时，即使假设空间很小也可能无法识别目标标签\scite{bendavid2010adaptation,bendavid2010impossibility}。这个问题与因果可识别性、时间依赖、在线动态比较及鲁棒优化的允许集合发生交叉。

估计源域与目标域的输入差异，只能说明可见分布相差多少，无法决定目标标签。若支持缺少覆盖或同一输入跨环境改变标签，无标签目标数据便可能同时符合两种要求相反预测的来源。迁移分析须先建立可识别联系，再计算有限样本代价；更精确的差异估计仍无法填补标签信息缺口。

多环境训练、持续变化与测试时适应进一步改变了数据关系。环境之间提供辨别信息时，多源观测可能区分偶然相关与稳定关系；在线更新可利用近期数据；交互收集还会随模型输出改变可见样本与评价分布。此时须说明环境生成、保持的机制及变化的可识别条件，不能只给固定分布风险附加一个误差项。

这些扩展共同保留了最初的预测问题：固定规则发展到候选类，样本保证加上计算要求，最坏容量补充算法选解，同分布风险推广到环境可识别性。后来的工具继续使用前面的条件，同时更细地规定数据来源、可见信息、比较基准与概率对象。接下来讨论的前沿问题也沿这些分别建立的关系展开。

\subsection{当前前沿}

图~\ref{fig:learning-theory-research-problems-map}从有限经验何时支持推广这一总问题出发，按问题与信息、泛化机制、实现与扩展分成三组，再列出八个相互制约的问题及部分代表课题。下面据此考察正在发展的方向，并指出已有条件尚不足以回答哪些问题。

\begin{figure*}[htbp]
  \centering
  \includegraphics[width=164.600mm]{figures/03-foundations/02-learning-theory/generated/learning-theory-map/figure.png}
  \caption{学习理论的问题与条件地图。顶层询问有限经验怎样支持可计算且可推广的学习；下方以蓝、绿、黄区分问题与信息、泛化机制、实现与扩展三组，这三组共列八个相互制约的问题。末端列表给出部分代表课题，不表示完整分类；连线表示讨论层次，不表示推导或因果关系。}
    \label{fig:learning-theory-research-problems-map}

\end{figure*}

\Needspace{20\baselineskip}
\paragraph{数据与算法相关的泛化解释}

当候选类的最坏容量远大于样本量，一致上界可能有效却宽到无法解释实际选解。研究于是考察间隔、范数、压缩、稳定性、信息依赖、参数后验与完整轨迹，以记录训练真正到达的结果。这些量本身也可能由训练数据选定，并随参数化和优化器改变；小数值只有连同选择过程的概率依赖一起受控，才支持泛化判断。

语言模型还需要先明确样本单位。文档、序列、词元与窗口通常有依赖，按独立样本累计词元可能高估有效信息。近期研究在明确损失与置信条件下，利用参数后验或数据分解得到严格低于对应平凡上界的泛化界\scite{lotfi2024llm,lotfi2024tokens}。它们仍依赖条件独立性、先验与数据选择。数值有用、条件可检验且能预测算法干预的风险界，仍需同时满足这些要求。

小泛化指标还可能只表现为相关。平坦性、频谱或稳定性与测试风险共同变化，并未自动解释改变初始化、样本顺序或时长的后果。把终点参数与训练动力学联系起来，需要理论对受控干预后的选解及风险作出可检验预测，而不只拟合已有结果。

\Needspace{10\baselineskip}
\paragraph{预训练、上下文学习与表示迁移}

冻结预训练表示减少了优化自由度，也可能遗漏目标信息；开放更多参数增加了可调整的自由度，可能改善任务适应，也可能增加下游依赖与遗忘。变化的方向和幅度须随表示、任务、数据和优化条件判断。因此预训练后的样本复杂度不能只按可训练参数计算，参数效率与统计效率也须分别衡量。分析需要量化预训练经验、表示选择及下游标签各提供多少任务信息。

上下文学习让参数不变的模型仍能根据提示示例调整预测，训练与使用便不再只由参数更新区分。受控理论可以将它分析为隐式贝叶斯推断，或把Transformer作为执行某类学习算法的对象\scite{xie2022bayesianicl,li2023icl}。这类解释均须给定预训练分布、上下文生成及任务族，不能由特定模型中的机制推断任意基础模型都遵循它。

上下文示例的顺序、相互依赖及新增示例引起的注意力或输出长度变化，还会改变评价条件。稳定性与非独立同分布分析尝试控制替换及依赖，长度泛化则研究训练长度之外的规则保持\scite{wang2026iclstability,huang2025length}。表示、算法模拟与序列外推因而相互联系。推进到真实系统时，预训练语料不透明且适应可能由多种机制共同产生，受控解释仍须检验其适用范围。

\Needspace{16\baselineskip}
\paragraph{统计、计算与可验证性的统一}

达到样本复杂度下界的阶，也不能保证训练、预测、存储或通信可承受。理论若调用优化预言机、独立采样器或无限精度运算，实现时就要补计有限迭代、随机误差、精度与停止条件。联合样本、时间、内存和通信预算的分析，正是将统计存在性落实到资源受限算法所需的步骤。

范数、覆盖常数、稳定性或分布差异等条件，若依赖未知总体，使用者就难以直接确认当前系统是否满足。将它们转成可检查证据，需要区分不可观测假设与可由日志、独立数据、重复运行或干预验证的条件，并核算验证本身的样本及资源成本。

\Needspace{10\baselineskip}
\paragraph{规模规律、能力变化与测试时计算}

大模型实验同时增加参数、数据与计算时，\term{缩放定律}（scaling law）以经验曲线记录资源与测试损失的关系，辅助预算分配\scite{kaplan2020scaling,hoffmann2022chinchilla}。曲线先描述已测范围，不能自动保证外推。数据重复、优化残差、模型形状或评价污染均可改变它；多种资源同步增长时，仅靠拟合也难分辨主导机制。

若要解释这种曲线，须区分近似、估计、优化与特征学习可能产生的相似幂律，再预测数据谱、训练时长或形状改变的后果。阈值化评价也可能制造能力突现，\term{顿悟}（grokking）则可能涉及特征形成、隐式偏差及测量尺度\scite{bordelon2025featurelaws,schaeffer2023emergence,power2022grokking}。能再现曲线尚不足以辨认机制，新的可区分实验预测才提供进一步证据。

推理中的采样、搜索与验证使计算预算从训练延伸到使用。某些任务和策略增加测试时计算可以改善结果，收益却仍由候选生成、选择器及预算共同决定\scite{snell2025testtime}。模型变大与同一模型搜索更久，需要分别分析对错误相关性、选择偏差及最终风险的影响。跨资源范围的外推还须保留这两类不同机制。

\Needspace{10\baselineskip}
\paragraph{变化环境中的持续学习保证}

预训练语料可能已含目标来源、相近样本或任务线索，使“未见环境”不能仅按微调数据界定。若忽略这些暴露，领域泛化评价可能把已有覆盖算成对新分布的外推\scite{mayilvahanan2025domaingeneralization}。因此应记录预训练来源、选模所见环境与最终测试环境，再说明哪些跨域关系能识别目标风险。

持续变化还需要将更新规则与到达数据的机制结合。被动数据流由外部环境继续生成；推荐和控制中的输出却可影响下一批可见数据。测试时适应的无标签代理目标可能随变化失效，更新流也可能被操纵。恢复复杂度研究达到规定性能所需的代价\scite{zhou2026ttalearnability}，变化下的共形方法研究协变量与后验同时漂移时的覆盖\scite{wang2025conformalshift}。变化范围、反馈与加权误差都是这些结论的条件，原交换性保证不能直接延用。

这些设定共同要求处理时间依赖及自适应交互。可用保证还需联合控制检测延迟、更新稳定性、历史任务退化与恢复时间，并规定停止或回退的条件。经典可学习性由此扩展到实际算法、预算和数据过程；条件能否检查，以及失效范围能否明确，也成为分析的一部分。

\section{本章小结}
\label{sec:statistical-framework-summary}

有限样本支持未来风险判断，首先取决于观测单位、比较与输出空间、损失、允许分布及随机来源。二分类中算法比较训练经验风险，目标却是总体期望风险。PAC用$\varepsilon$限定风险或超额风险，用$\delta$限定随机训练失败概率。可实现分布有零风险目标，不可知分布允许噪声与表达限制，改以类内下确界为基准。两种设定都要求同一充分样本量覆盖全部允许分布。

低经验风险不能独自完成这项推断。自由标签空间中的没有免费的午餐下界来自未见标签信息不足，更多计算无法消除。加入假设限制可以建立输入间的联系，却可能排除所需规则。贝叶斯风险、近似误差与类内超额风险分别定位输入、空间与算法选择的贡献；经验优化误差和泛化差则比较不同的量，须保留各自基准。

后续理论在这些关系上改变分析对象：候选容量之外考察稳定性、信息依赖与网络选解，同分布之外考察环境联系与可识别性。标签结构、损失、参数化和数据来源的变化，都要重新核对表示、有限样本与计算条件。规模规律和持续适应仍受这些基础条件约束。

结构风险最小化已经给出一种连接：同时控制各空间的经验偏差，算法即可用拟合加惩罚与最佳参照空间比较。下一章将为二分类建立提供这种偏差控制的容量条件。统计样本界成立后，还须另给能停机并满足时间预算的求解程序，才能获得可计算或高效学习。

\paragraph{延伸阅读}
Shalev-Shwartz与Ben-David的\textit{Understanding Machine Learning}第2--5章展开PAC、ERM与没有免费的午餐，第7章讨论SRM及描述长度。阅读这些结果时，可核对各自比较基准及算法先于分布的量词顺序，从而判断结论适用于哪一类学习问题\scite{shalevshwartz2014}。

Delle Rose等的可计算PAC研究进一步区分统计存在与有效见证\scite{dellerose2023}。知道某种预测不可能出现，尚未给出有限时间找出它的程序。第~\ref{sec:computable-binary-learning}节会以定义和例子落实这种区别，可与原文的表示及输出约定对照阅读。
